\documentclass[10pt]{beamer}
\usetheme{Madrid}
\usecolortheme{whale}
\usepackage[utf8]{inputenc}
\usepackage[T1]{fontenc}
\usepackage[spanish,es-tabla,es-nodecimaldot]{babel}
\usepackage{amsmath,amssymb}
\usepackage{booktabs}
\usepackage{array}
\newcolumntype{L}[1]{>{\raggedright\arraybackslash}p{#1}}
\usepackage{tikz}
\usetikzlibrary{babel}

\setbeamertemplate{navigation symbols}{}
\newcommand{\dato}[1]{\textbf{\color{blue!60!black}#1}}

\title[Laguna de Aculeo: resultados]{Superficie hídrica de la Laguna de Aculeo}
\subtitle{Resultados de los modelos y diferencias con la Laguna Torca}
\author[J. P. Maureira]{Juan Pablo Maureira}
\institute[UTalca]{Magíster en Data Science\\Universidad de Talca}
\date{\today}

\begin{document}

\begin{frame}
\titlepage
\end{frame}

% ------------------------------------------------------------
\begin{frame}{Datos y cambios respecto de Torca}
\begin{itemize}
\item \textbf{Superficie:} 687 escenas Landsat 5, 7, 8 y 9 (1985--2024), NDWI en Google Earth Engine. Se usan solo escenas con al menos 90\,\% del polígono visible.
\item \textbf{Clima:} CR2MET (1979--2019) para la regresión, y ERA5-Land (hasta 2024) para el balance hídrico.
\item \textbf{Proyecciones:} CR2, RCP~2.6 y RCP~8.5, hasta 2095.
\end{itemize}
\vspace{2mm}
\textbf{Dos decisiones distintas a Torca}, porque aquí la sequía es real:
\begin{itemize}
\item Las áreas bajas \textbf{no se borran} como nube.
\item El cambio de sensor se calibra con \textbf{pares de escenas casi simultáneas}, en vez de restar la media de cada época (eso habría borrado la sequía).
\end{itemize}
\vspace{2mm}
\centering
\small
\begin{tabular}{lrrl}
\toprule
Par de sensores & Pares & Factor (IC 95\,\%) & Decisión \\
\midrule
Landsat 5 $\to$ 7 & 83 & 0.99 (0.99--0.99) & Se aplica \\
Landsat 7 $\to$ 8 & 30 & 1.04 (1.03--1.07) & Se aplica \\
Landsat 9 $\to$ 8 & 16 & 0.99 (0.95--1.03) & No se aplica \\
\bottomrule
\end{tabular}
\end{frame}

% ------------------------------------------------------------
\begin{frame}{Aculeo se secó por completo y volvió a llenarse}
\begin{columns}[T]
\begin{column}{0.55\textwidth}
\centering
\begin{tikzpicture}[x=0.68cm,y=0.95cm,font=\scriptsize]
\draw[gray!50] (0,0) rectangle (8,3.5);
\foreach \v/\l in {0/0,0.808/300,1.615/600,2.423/900,3.231/1200} {\draw[gray!30] (0,\v)--(8,\v); \node[left] at (0,\v) {\l};}
\foreach \x/\l in {0/1985,2/1995,4/2005,6/2015,8/2025} {\node[below] at (\x,0) {\l};}
\draw[teal!70!black,very thick] plot coordinates {(0.20,2.825) (0.60,3.338) (0.80,3.101) (1.00,3.047) (1.20,2.692) (1.40,3.004) (1.60,2.897) (1.80,2.841) (2.00,2.731) (2.20,3.034) (2.40,3.197) (2.60,3.295) (2.80,2.680) (3.00,3.123) (3.20,3.181) (3.40,2.917) (3.60,3.096) (3.80,3.298) (4.00,2.711) (4.20,3.131) (4.40,3.122) (4.60,2.790) (4.80,3.052) (5.00,2.592) (5.20,2.895) (5.40,2.644) (5.60,2.703) (5.80,1.992) (6.00,1.043) (6.20,0.380) (6.40,0.546) (6.60,0.115) (6.80,0.002) (7.00,0.004) (7.20,0.007) (7.40,0.000) (7.60,0.446) (7.80,1.747)};
\draw[orange!80!black,thick,dashed] plot coordinates {(0.20,0.348) (0.40,0.336) (0.60,0.352) (0.80,0.348) (1.00,0.339) (1.20,0.342) (1.80,0.320) (2.00,0.338) (2.20,0.330) (2.40,0.342) (2.60,0.313) (2.80,0.325) (3.00,0.318) (3.20,0.328) (3.40,0.332) (3.60,0.328) (3.80,0.335) (4.00,0.341) (4.20,0.342) (4.40,0.338) (4.60,0.315) (4.80,0.335) (5.00,0.327) (5.20,0.338) (5.60,0.359) (5.80,0.360) (6.00,0.344) (6.20,0.358) (6.40,0.362) (6.60,0.369) (6.80,0.350) (7.00,0.365) (7.20,0.381) (7.40,0.368) (7.60,0.381) (7.80,0.379)};
\node[rotate=90] at (-1.3,1.75) {Área media anual (ha)};
\node[teal!70!black] at (2.6,2.3) {Aculeo};
\node[orange!80!black] at (2.6,0.65) {Torca (misma escala)};
\end{tikzpicture}
\end{column}
\begin{column}{0.43\textwidth}
\small
\begin{itemize}
\item 1985--2013: \dato{963 a 1\,240 ha} por año (media ${\approx}1\,100$ ha).
\item 2014--2018: colapso, de 740 a 43 ha.
\item 2019--2022: \dato{seca} (menos de 3 ha).
\item Septiembre 2023: ${\approx}460$ ha tras un solo invierno.
\item Fines de 2024: \dato{${\approx}1\,060$ ha}.
\end{itemize}
\end{column}
\end{columns}
\vspace{2mm}
\small Es una señal unas 100 veces más grande que la de la Torca, con un colapso y dos rellenados reales.
\end{frame}

% ------------------------------------------------------------
\begin{frame}{Correlación con el clima: responde a la lluvia de varios años}
\centering
\footnotesize
\begin{tabular}{lrrrr}
\toprule
Variable & $\rho$ sin ciclo & $R^2$ & $\rho$ 1985--2012 & $\rho$ 2013--2019 \\
\midrule
Lluvia del mes & $-0.07$ & 0.5\,\% & $-0.16$ & $-0.07$ \\
Lluvia acumulada 12 meses & 0.30 & 9\,\% & 0.24 & $-0.20$ \\
Lluvia acumulada 24 meses & 0.40 & 16\,\% & 0.25 & $-0.35$ \\
Lluvia acumulada 36 meses & \dato{0.47} & \dato{22\,\%} & 0.20 & $-0.56$ \\
Temperatura máxima & $-0.06$ & 0.4\,\% & 0.06 & $-0.28$ \\
\bottomrule
\end{tabular}
\vspace{3mm}
\begin{itemize}
\item A diferencia de Torca, aquí sí aparece una asociación con la \textbf{lluvia acumulada de 2 a 3 años}.
\item Pero el tamaño de muestra efectivo es chico (${\approx}12$ meses independientes para el acumulado de 36 meses) por la autocorrelación.
\item Y el signo \textbf{se invierte después de 2013}: la laguna siguió secándose en 2016--2018 aunque la lluvia ya se había recuperado.
\end{itemize}
\end{frame}

% ------------------------------------------------------------
\begin{frame}{Modelos con validación temporal: gana la persistencia}
\centering
\small
\begin{tabular}{lrr}
\toprule
Modelo & RMSE (ha) & Ganancia* \\
\midrule
Persistencia: último valor (referencia) & \dato{287} & 47\,\% \\
Random Forest con clima & 498 & 8\,\% \\
XGBoost con clima & 512 & 6\,\% \\
Ridge con clima & 514 & 5\,\% \\
Random Forest sin clima & 542 & 0\,\% \\
Ridge sin clima & 542 & 0\,\% \\
Climatología mensual (referencia) & 542 & -- \\
Media histórica (referencia) & 551 & $-2$\,\% \\
\bottomrule
\end{tabular}\\[1mm]
{\scriptsize *Reducción del RMSE respecto de la climatología mensual.}
\vspace{3mm}
\begin{itemize}
\item El clima \textbf{sí aporta}: baja el error entre 28 y 44 ha en los tres modelos (en Torca no aportaba nada).
\item Pero ningún modelo con clima se acerca a la persistencia: su error (${\approx}500$ ha) es casi la mitad del tamaño de la laguna.
\item La regresión con lluvia acumulada no puede aprender una laguna que resiste años secos y de pronto colapsa.
\end{itemize}
\end{frame}

% ------------------------------------------------------------
\begin{frame}{Balance hídrico con ERA5-Land (2020--2024)}
\begin{columns}[T]
\begin{column}{0.48\textwidth}
\small
Un ``balde'': entra la lluvia sobre un umbral, sale una pérdida constante y tiene un tope. Calibrado en 1985--2019 y probado en 2020--2024 sin reajustar.
\vspace{2mm}

\centering
\begin{tabular}{lr}
\toprule
Referencia (2020--2024) & RMSE (ha) \\
\midrule
Media histórica & 805 \\
Persistencia (0 ha) & 332 \\
Balance hídrico & \dato{259} \\
\bottomrule
\end{tabular}
\end{column}
\begin{column}{0.5\textwidth}
\small
\centering
\begin{tabular}{lrr}
\toprule
Invierno & Lluvia may--sep & Área oct--dic \\
\midrule
2017 & 517 mm & 328 ha \\
2020 & 393 mm & \dato{0 ha} \\
2023 & 337 mm & \dato{507 ha} \\
2024 & 546 mm & 1\,049 ha \\
\bottomrule
\end{tabular}
\end{column}
\end{columns}
\vspace{3mm}
\begin{itemize}
\item Mejora a la persistencia, pero \textbf{no reproduce el rellenado de 2023} (predice 0 ha; se observaron 460--530 ha) y el coeficiente de evaporación quedó en cero.
\item En 2020 llovió más que en 2023 y la laguna no se llenó: la lluvia mensual en una celda de ${\approx}9$ km no explica la respuesta.
\end{itemize}
\end{frame}

% ------------------------------------------------------------
\begin{frame}{Proyecciones RCP: el escenario sí importa, pero el modelo no es confiable}
\centering
\footnotesize
\begin{tabular}{lrrrr}
\toprule
Horizonte & RCP~2.6 (ha) & RCP~8.5 (ha) & 8.5 $-$ 2.6 (ha) & RCP~8.5 $-$ control (ha) \\
\midrule
2025--2040 & 874 & 827 & $-47$ & $-260$ \\
2041--2070 & 872 & 774 & $-98$ & $-314$ \\
2071--2095 & 917 & 718 & \dato{$-199$} & $-369$ \\
\bottomrule
\end{tabular}
\vspace{3mm}
\begin{itemize}
\item A diferencia de Torca, el modelo proyecta \textbf{menos agua con RCP~8.5}, en la dirección físicamente esperada.
\item La diferencia crece con el horizonte, hasta ${\approx}200$ ha en 2071--2095.
\item Pero el modelo tiene un error de ${\approx}500$ ha y pierde contra la persistencia: estas cifras son un \textbf{análisis de sensibilidad}, no una predicción.
\item Además usa las proyecciones CR2 descargadas para Torca, no para la cuenca de Aculeo.
\end{itemize}
\end{frame}

% ------------------------------------------------------------
\begin{frame}{Diferencias entre Aculeo y Torca}
\centering
\footnotesize
\setlength{\tabcolsep}{4pt}
\renewcommand{\arraystretch}{1.15}
\begin{tabular}{L{2.6cm}L{4.0cm}L{4.0cm}}
\toprule
 & \textbf{Laguna Torca} & \textbf{Laguna de Aculeo} \\
\midrule
Rango del área anual & 116--142 ha (${\approx}\pm10$\,\%) & 0--1\,240 ha (se secó) \\
Efecto del sensor & Salto de 11 ha, confundido con la época & 1--4\,\%, calibrado con pares de escenas \\
Mayor $|\rho|$ con el clima & 0.21 mensual; 0.27 anual (no signif.) & 0.47 (lluvia acumulada 36 meses) \\
Escala de respuesta & Ninguna clara & Varios años de lluvia \\
Aporte del clima al RMSE & Nulo o negativo & Baja 28--44 ha \\
Mejor modelo & Empate: Ridge estacional y persistencia (${\approx}7.7$ ha) & Persistencia (287 ha) \\
RCP~8.5 $-$ RCP~2.6, 2071--2095 & ${\approx}0$ ha (MICE: $+22.6$ ha, artefacto) & $-199$ ha (sensibilidad) \\
Dirección del efecto & Sin efecto & Menos agua con más sequía \\
Qué falta explicar & Casi todo: sistema lacustre, manejo, napa & Rellenado 2023 y colapso 2016--2018: extracción, napa, intensidad de tormentas \\
\bottomrule
\end{tabular}
\end{frame}

% ------------------------------------------------------------
\begin{frame}{Conclusiones}
\begin{itemize}
\item Aculeo tiene una señal enorme y limpia (sesgo entre sensores de 1--4\,\%): es mejor laguna de estudio que Torca.
\item El clima \textbf{sí deja huella}: la lluvia acumulada de 3 años explica ${\approx}22$\,\% de la variación y mejora los modelos, cosa que en Torca no ocurre.
\item Aun así, ningún modelo con clima supera a la persistencia, y el balance hídrico no reproduce el rellenado de 2023.
\item La proyección va en la dirección esperada (hasta $-199$ ha con RCP~8.5), pero con un error de ${\approx}500$ ha solo vale como sensibilidad.
\item En ambas lagunas, lluvia y temperatura mensuales \textbf{no bastan} para proyectar a 100 años. En Aculeo faltan datos de extracción, napas y lluvia diaria; en Torca, de la conexión con su sistema lacustre.
\end{itemize}
\end{frame}

\end{document}
